\section*{Bab \ref{ch:b2:structures} --- Himpunan dan Struktur}

\begin{solution}{exo:b2:structures:1}
\emph{Himpunan bagian hingga dari $\N$:} \omterm{def:b2:structures:countable}{terbilang} --- himpunan semua
himpunan bagian $\intint{0}{n}$ adalah hingga, dan himpunan bagian yang
hingga membentuk gabungan \omterm{def:b2:structures:countable}{terbilang} atas $n$ dari himpunan-himpunan itu
(\cref{prop:b2:structures:countablestable}~(3)); tak hingga karena ia
memuat semua himpunan tunggal.

\emph{Semua himpunan bagian $\N$:} tidak \omterm{def:b2:structures:countable}{terbilang}, menurut teorema
Cantor (\cref{thm:b2:structures:cantor}~(1) dengan $E = \N$).

\emph{$\R \setminus \Q$:} tidak \omterm{def:b2:structures:countable}{terbilang} --- sebab jika tidak, $\R = \Q
\cup (\R\setminus\Q)$ akan menjadi gabungan dua \omterm{def:b2:structures:countable}{himpunan terbilang}, yang
bertentangan dengan \cref{thm:b2:structures:cantor}~(2).

\emph{Polinomial atas $\Q$:} \omterm{def:b2:structures:countable}{terbilang} --- polinomial berderajat
$\leq n$ terinjeksi ke $\Q^{n+1}$ (hasil kali hingga atas \omterm{def:b2:structures:countable}{himpunan
terbilang}), lalu ambil gabungan atas $n$.

\emph{Barisan biner yang akhirnya nol:} \omterm{def:b2:structures:countable}{terbilang} --- barisan itu
berpadanan bijektif dengan himpunan bagian hingga $\N$ (lewat penyangganya).
\end{solution}

\begin{solution}{exo:b2:structures:2}
Hitung unsur demi unsur, dengan faktor kanan dikerjakan lebih dulu.
$\sigma\tau$ mengirim $1 \mapsto \sigma(1) = 4$, $\;2 \mapsto \sigma(3)
= 5$, $\;3 \mapsto \sigma(7) = 7$, $\;4 \mapsto \sigma(4) = 2$, $\;5
\mapsto \sigma(5) = 3$, $\;6 \mapsto \sigma(6) = 1$, $\;7 \mapsto
\sigma(2) = 6$:
\[
\sigma\tau = (1\,4\,2\,5\,3\,7\,6),
\]
sebuah \omterm{def:b2:structures:sn}{siklus} berpanjang $7$. Serupa itu $\tau\sigma$ mengirim $1
\mapsto \tau(4) = 4$, $\;2 \mapsto \tau(6) = 6$, $\;3 \mapsto \tau(5) =
5$, $\;4 \mapsto \tau(2) = 3$, $\;5 \mapsto \tau(3) = 7$, $\;6 \mapsto
\tau(1) = 1$, $\;7 \mapsto \tau(7) = 2$:
\[
\tau\sigma = (1\,4\,3\,5\,7\,2\,6),
\]
juga \omterm{def:b2:structures:sn}{siklus} berpanjang $7$ (sesuai dugaan: $\sigma\tau$ dan $\tau\sigma$
sekonjugasi, jadi tipe siklusnya sama).

\omterm{def:b2:structures:generated}{Orde} dan tanda: $\sigma$ bertipe \omterm{def:b2:structures:sn}{siklus} $(4,2)$, jadi berorde
$\operatorname{lcm}(4,2) = 4$ dan bertanda $(-1)^3(-1)^1 = +1$; $\tau$
adalah \omterm{def:b2:structures:sn}{siklus} berpanjang $3$, jadi berorde $3$ dan bertanda $+1$; kedua
hasil kali tadi \omterm{def:b2:structures:sn}{siklus} berpanjang $7$, jadi berorde $7$ dan bertanda
$(-1)^6 = +1$.

$\sigma^{2026}$: karena $2026 = 4 \times 506 + 2$, maka $\sigma^{2026} =
\sigma^2 = (1\,2)(4\,6)$ (kuadratkan \omterm{def:b2:structures:sn}{siklus} berpanjang $4$; \omterm{def:b2:structures:sn}{transposisinya}
lenyap ketika dikuadratkan).
\end{solution}

\begin{solution}{exo:b2:structures:3}
$\{0,1\}^{\N} \to \mathcal{P}(\N)$: sebuah barisan dipetakan ke
penyangganya --- sebuah bijeksi (fungsi indikator), tanpa perlu teorema.

$\{0,1\}^{\N} \to \intcc{0}{1}$: pemetaan basis $3$, yakni $(a_n) \mapsto
\sum 2a_n 3^{-n-1}$, bersifat injektif (dua barisan yang berbeda mulai
berselisih pada peringkat $N$; ekornya tak sanggup menutup jurang sebesar
$2\cdot 3^{-N-1}$, sebab $\sum_{n > N} 2\cdot 3^{-n-1} = 3^{-N-1} <
2\cdot3^{-N-1}$).

$\intcc{0}{1} \to \{0,1\}^{\N}$: uraian biner, dengan memilih (misalnya)
uraian yang tidak berakhir dengan $1$ berulang: injektif.

Menurut Cantor--Bernstein (\cref{thm:b2:structures:cantorbernstein})
yang diterapkan pada dua injeksi terakhir, $\intcc{0}{1}$ dan
$\{0,1\}^{\N}$ \omterm{def:b2:structures:countable}{ekuipoten}, jadi ketiga himpunan itu \omterm{def:b2:structures:countable}{ekuipoten}.
\end{solution}

\begin{solution}{exo:b2:structures:4}
Misalkan $c = ab = ba$ dan $d = \operatorname{ord}(c)$. Pertama, $c^{mn} =
a^{mn} b^{mn} = e$ (kekomutatifan memungkinkan pemisahan pangkat), jadi
$d \mid mn$. Sebaliknya $c^d = e$ memberi $a^d = b^{-d}$; unsur ini
berada di $\langle a\rangle \cap \langle b\rangle$, yaitu subgrup yang
\omterm{def:b2:structures:generated}{ordenya} membagi $m$ sekaligus $n$ (Lagrange pada masing-masing \omterm{def:b2:structures:generated}{grup
siklik}), sehingga subgrup itu trivial: $a^d = b^d = e$, jadi $m \mid d$
dan $n \mid d$, dan karena saling prima $mn \mid d$. Maka $d = mn$.

Di $\mathfrak{S}_3$: ambil $a = (1\,2)$ (berorde $2$) dan $b =
(1\,2\,3)$ (berorde $3$), yang \omterm{def:b2:structures:generated}{ordenya} saling prima tetapi tidak
komutatif: $ab = (2\,3)$ berorde $2 \neq 6$ --- memang $\mathfrak{S}_3$
tidak punya unsur berorde $6$. Kekomutatifan memang penting.
\end{solution}

\begin{solution}{exo:b2:structures:5}
Pasangkan setiap $x \in G$ dengan $x^{-1}$. Pasangan $\{x, x^{-1}\}$
dengan $x \neq x^{-1}$ beranggotakan dua unsur dan memilah gabungannya;
unsur yang tersisa persis yang memenuhi $x = x^{-1}$, yakni $x^2 =
e$. Karena $\abs G$ genap dan pasangan berunsur dua itu meliputi
sejumlah genap unsur, himpunan $\{x : x^2 = e\}$ berkardinalitas genap;
ia memuat $e$, jadi ia memuat sedikitnya satu unsur lain $x \neq
e$ --- yakni sebuah unsur berorde $2$.
\end{solution}

\begin{solution}{exo:b2:structures:6}
Setiap unsur $A_n$ adalah hasil kali sejumlah genap \omterm{def:b2:structures:sn}{transposisi}
(\cref{thm:b2:structures:signature}: uraikan atas \omterm{def:b2:structures:sn}{transposisi};
banyaknya genap karena tandanya $+1$). Cukuplah menuliskan setiap hasil
kali dua \omterm{def:b2:structures:sn}{transposisi} memakai \omterm{def:b2:structures:sn}{siklus} berpanjang $3$:
\[
(a\,b)(a\,c) = (a\,c\,b),
\qquad
(a\,b)(c\,d) = (a\,c\,b)(a\,c\,d) \quad (\text{dengan } a,b,c,d \text{ berbeda}),
\]
(periksa lewat penilaian), dan $(a\,b)(a\,b) = \mathrm{id}$. Jadi
\omterm{def:b2:structures:sn}{siklus} berpanjang $3$ membangun $A_n$.
\end{solution}

\begin{solution}{exo:b2:structures:7}
\emph{$(\Q,+) \to (\Z,+)$:} hanya morfisma nol. Untuk sembarang $x$ dan
setiap $n \geq 1$, $f(x) = n f\bigl(\frac xn\bigr)$ habis dibagi
$n$ di $\Z$; satu-satunya bilangan bulat yang habis dibagi setiap $n$
adalah $0$, jadi $f(x) = 0$ untuk setiap $x$.

\emph{$(\Z/n\Z, +) \to (\Z/m\Z, +)$:} sebuah morfisma ditentukan oleh
$c = f(\overline 1)$, yang wajib memenuhi $n c \equiv 0 \pmod m$,
yakni $c$ merupakan kelipatan $\frac{m}{\gcd(m,n)}$; ada
$\gcd(m,n)$ kelas semacam itu, dan tiap pilihan memang mendefinisikan
morfisma (faktorkan $k \mapsto kc$ lewat $\Z/n\Z$ dengan sifat universal).

\emph{$(\Q, +) \to (\Q_+^*, \times)$:} hanya morfisma trivial. Jika
$f(x) = y$, maka untuk setiap $n$, $y = f(n \cdot \frac xn) =
f(\frac xn)^n$ merupakan pangkat ke-$n$ di $\Q_+^*$. Namun bilangan
rasional $y \neq 1$ tak mungkin menjadi pangkat ke-$n$ untuk setiap $n$:
ada bilangan prima yang muncul di $y$ dengan pangkat $v$ tak nol, dan $n
\nmid v$ untuk $n > \abs v$ (pangkat pada suatu pangkat ke-$n$ selalu
kelipatan $n$, menurut ketunggalan pemfaktoran). Jadi $f \equiv 1$.
\end{solution}

\begin{solution}{exo:b2:structures:8}
$360 = 2^3 \cdot 3^2 \cdot 5$, jadi
$\varphi(360) = 360\bigl(1 - \tfrac12\bigr)\bigl(1 -
\tfrac13\bigr)\bigl(1 - \tfrac15\bigr) = 360 \cdot \tfrac12 \cdot
\tfrac23 \cdot \tfrac45 = 96$.

Sistemnya: modulus $8, 9, 5$ saling prima sepasang demi sepasang, dengan
hasil kali $360$. Dari $x \equiv 3 \pmod 8$ dan $x \equiv 5 \pmod 9$:
$x = 3 + 8k$ dengan $3 + 8k \equiv 5 \pmod 9$, yakni $-k \equiv 2$, $k
\equiv -2 \equiv 7 \pmod 9$, sehingga $x \equiv 3 + 56 = 59 \pmod{72}$.
Lalu $59 + 72\ell \equiv 2 \pmod 5$: $4 + 2\ell \equiv 2$, $2\ell \equiv
3 \equiv 8$, $\ell \equiv 4 \pmod 5$, jadi $x \equiv 59 + 288 = 347 \pmod{360}$.

Dua angka terakhir $3^{2026}$: modulo $4$, $3^{2026} = 9^{1013} \equiv
1$. Modulo $25$: $\varphi(25) = 20$ dan $2026 = 20\cdot101 + 6$, jadi
$3^{2026} \equiv 3^6 = 729 \equiv 4 \pmod{25}$. Selesaikan $x \equiv 1
\pmod 4$, $x \equiv 4 \pmod{25}$: dari $x = 4 + 25k \equiv 1 \pmod 4$
diperoleh $k \equiv 1 \pmod 4$, jadi $x \equiv 29 \pmod{100}$. Dua angka
terakhirnya adalah $29$.
\end{solution}

\begin{solution}{exo:b2:structures:9}
Misalkan $A$ daerah integral yang hingga dan $a \in A$, $a \neq 0$.
Pemetaan $x \mapsto ax$ bersifat injektif ($ax = ay \implies a(x - y) = 0
\implies x = y$, sebab tidak ada pembagi nol); pemetaan injektif dari
himpunan hingga ke dirinya sendiri pastilah surjektif (jilid Tahun ke-1,
kesetaraan sarang merpati). Jadi $1 = ab$ untuk suatu $b$: setiap unsur
tak nol punya invers, sehingga $A$ adalah lapangan.

$\Z/n\Z$: jika $n$ prima, ia daerah integral ($n \mid ab
\implies n \mid a$ atau $n \mid b$, lewat lema Euklides), hingga, jadi
lapangan; jika $n = rs$ komposit, maka $\overline r\,\overline s =
\overline 0$ memperlihatkan adanya pembagi nol.
\end{solution}

\begin{solution}{exo:b2:structures:10}
Misalkan $m = \max\{\operatorname{ord}(x) : x \in G\}$, yang dicapai di $a$.

\emph{Klaim: \omterm{def:b2:structures:generated}{orde} setiap $x \in G$ membagi $m$.} Andaikan ada
$x$ berorde $q$ dengan $q \nmid m$: maka ada pangkat prima $p^k$ yang
membagi $q$ tetapi tidak membagi $m$. Tulis $m = p^j m'$ dengan $p \nmid
m'$ dan $j < k$. Unsur $a^{p^j}$ berorde $m'$; unsur
$x^{q/p^k}$ berorde $p^k$; kedua \omterm{def:b2:structures:generated}{orde} itu saling prima dan kedua
unsurnya komutatif ($G \subseteq K^*$ abelian), jadi menurut
\cref{exo:b2:structures:4} hasil kalinya berorde $p^k m' > p^j m'
= m$: bertentangan dengan kemaksimalan.

Jadi setiap $x \in G$ memenuhi $x^m = 1$: polinomial $X^m - 1$ punya
sedikitnya $\abs G$ akar di lapangan $K$, sehingga $\abs G \leq m$
(polinomial tak nol berderajat $m$ punya paling banyak $m$ akar, jilid
Tahun ke-1). Namun $m = \operatorname{ord}(a) \leq \abs G$ menurut
Lagrange. Maka $m = \abs G$ dan $\langle a \rangle$, yang berkardinalitas
$m = \abs G$, adalah seluruh $G$: \omterm{def:b2:structures:generated}{siklik}.

Untuk $K = \Z/p\Z$: $(\Z/p\Z)^*$ adalah subgrup hingga dari $K^*$, jadi
\omterm{def:b2:structures:generated}{siklik} (berorde $p - 1$).
\end{solution}

\begin{solution}{exo:b2:structures:11}
\emph{Tidak \omterm{def:b2:structures:generated}{siklik}:} subgrup $\langle \frac pq\rangle$ terdiri atas
kelipatan bulat $\frac pq$, yang semuanya berpenyebut pembagi $q$ (dalam
bentuk paling sederhana); karena itu subgrup itu melewatkan
$\frac{1}{2q}$. Tidak ada pembangun tunggal yang mampu menjangkau
penyebut $\Q$ yang tak terbatas.

\emph{Tidak \omterm{def:b2:structures:generated}{dibangun} oleh berhingga unsur:} subgrup yang \omterm{def:b2:structures:generated}{dibangun} oleh
$\frac{p_1}{q_1}, \dots, \frac{p_k}{q_k}$ terdiri atas bilangan rasional
yang penyebutnya membagi $Q = q_1 \cdots q_k$ (kombinasi bulat
berpenyebut pembagi $Q$): subgrup itu melewatkan $\frac{1}{2Q}$.

\emph{Subgrup yang \omterm{def:b2:structures:generated}{dibangun} oleh berhingga unsur adalah \omterm{def:b2:structures:generated}{siklik}:} dengan
$Q$ seperti di atas, subgrup $H = \langle \frac{p_1}{q_1}, \dots,
\frac{p_k}{q_k}\rangle$ termuat di $\frac{1}{Q}\Z$. Pemetaan $x
\mapsto Qx$ adalah isomorfisma dari $\frac1Q\Z$ pada $\Z$ yang membawa
$H$ ke suatu subgrup $\Z$, yaitu $n\Z$ untuk suatu $n$ (jilid Tahun
ke-1): jadi $H = \frac{n}{Q}\Z$ \omterm{def:b2:structures:generated}{siklik}, \omterm{def:b2:structures:generated}{dibangun} oleh
$\frac nQ$.
\end{solution}

\begin{solution}{exo:b2:structures:12}
\emph{Sebuah himpunan bagian \omterm{def:b2:structures:countable}{terbilang}.} Misalkan $E$ tak hingga.
Bangun $a_0, a_1, a_2, \dots$ secara induktif: $E$ tak kosong, jadi
pilih $a_0 \in E$; jika $a_0, \dots, a_n$ sudah terpilih, maka $E
\setminus \{a_0, \dots, a_n\}$ tak kosong ($E$ tidak hingga), jadi pilih
$a_{n+1}$ di sana. Semua $a_n$ berbeda sepasang demi sepasang menurut
konstruksinya, jadi $A = \{a_n : n \in \N\}$ himpunan bagian $E$ yang \omterm{def:b2:structures:countable}{terbilang}.

\emph{Tak hingga $\implies$ \omterm{def:b2:structures:countable}{ekuipoten} dengan himpunan bagian sejati.}
Tetapkan $f \colon E \to E \setminus \{a_0\}$ oleh $f(a_n) = a_{n+1}$
dan $f(x) = x$ untuk $x \notin A$. Pemetaan itu injektif (kedua
potongannya injektif dengan peta yang saling lepas) dan surjektif pada
$E \setminus \{a_0\}$: setiap $a_{n+1}$ tercapai, setiap $x \notin A$
tercapai. Jadi $E$ \omterm{def:b2:structures:countable}{ekuipoten} dengan himpunan bagian sejati $E \setminus \{a_0\}$.

\emph{Sebaliknya.} Jika $E$ hingga dan $g \colon E \to F$ adalah
bijeksi pada $F \subseteq E$ dengan $F \neq E$, maka $g$ merupakan
injeksi $E$ ke dirinya sendiri yang tidak surjektif, dan itu
bertentangan dengan asas sarang merpati (jilid Tahun ke-1: pemetaan
injektif dari himpunan hingga ke dirinya sendiri pasti bijektif). Jadi
himpunan yang \omterm{def:b2:structures:countable}{ekuipoten} dengan himpunan bagian sejatinya pastilah tak hingga.
\end{solution}

\begin{solution}{pb:b2:structures:1}
\textbf{1.} Sebuah konfigurasi memberikan kepada masing-masing dari $16$
sel satu dari $16$ isi (ubin $1$--$15$ atau kekosongan $b = 16$),
masing-masing tepat sekali: itu persis sebuah bijeksi $\intint1{16} \to
\intint1{16}$, yakni unsur $\mathfrak{S}_{16}$; banyaknya $16!
= 20\,922\,789\,888\,000$. Satu langkah sah menggeser satu ubin yang
bertetangga dengan kekosongan, jadi banyaknya langkah sama dengan
banyaknya tetangga sel kekosongan: $2$ untuk keempat sel pojok,
$3$ untuk kedelapan sel tepi, dan $4$ untuk keempat sel dalam.

\textbf{2.} Setelah geseran, sel $p$ memuat isi lama
$c$ dan sel $c$ memuat kekosongan; semua sel lain tak tersentuh:
$\sigma'(p) = \sigma(c)$, $\sigma'(c) = \sigma(p) = 16$, dan
$\sigma' = \sigma$ pada sel lainnya. Itu persis $\sigma' = \sigma
\circ (p\ c)$. Karena $\varepsilon$ adalah morfisma dan
$\varepsilon\bigl((p\ c)\bigr) = -1$, kita peroleh
$\varepsilon(\sigma') = -\varepsilon(\sigma)$.

\textbf{3.} Sel bertetangga berselisih satu langkah tepat pada salah
satu dari kedua koordinatnya, jadi paritas $i + j$ berubah: $\chi$
bernilai berlawanan pada sel yang bertetangga. Satu langkah memindahkan
kekosongan dari $p$ ke $c$ yang bertetangga, sehingga membalik
$\chi(\text{sel kekosongan})$. Sepanjang jalan tertutup yang ditempuh
kekosongan, $\chi$ terbalik sekali tiap langkah lalu kembali ke nilai
awalnya: jadi banyaknya langkah genap.

\textbf{4.} Menurut pertanyaan 2 dan 3, satu langkah membalik kedua
faktor $I(\sigma) = \varepsilon(\sigma)\chi(\sigma^{-1}(16))$, sehingga
hasil kalinya tak berubah. Untuk konfigurasi tersusun:
$\varepsilon(\mathrm{id}) = +1$ dan kekosongan berada di sel $16$,
baris $4$, kolom $4$, jadi $\chi(16) = (-1)^{8} = +1$, sehingga
$I(\mathrm{id}) = +1$.

\textbf{5.} $\sigma_L$ adalah \omterm{def:b2:structures:sn}{transposisi} $(14\ 15)$ atas sel:
$\varepsilon(\sigma_L) = -1$; kekosongannya di rumah, jadi $\chi(16) =
+1$ dan $I(\sigma_L) = -1 \neq +1 = I(\mathrm{id})$. Karena $I$
terpelihara oleh setiap langkah, tidak ada rangkaian langkah yang
menghubungkan $\sigma_L$ dengan $\mathrm{id}$. Hadiahnya aman secara struktural.

\textbf{6.} Tetapkan sebuah sel $p$ dan dua sel lain $c \neq d$ yang
berbeda dari $p$, lalu ambil $\tau_0 = (c\ d)$. Pada himpunan
konfigurasi yang kekosongannya di $p$, pemetaan $\sigma \mapsto \sigma
\circ \tau_0$ adalah involusi (ia memelihara $\sigma(p) = 16$
karena $\tau_0$ membiarkan $p$ tetap) dan membalik $\varepsilon$,
sehingga membalik $I$: pemetaan itu memasangkan konfigurasi ber-$I = +1$
secara bijektif dengan konfigurasi ber-$I = -1$. Jadi tiap satu dari $16$
kedudukan kekosongan menyumbang $15!/2$ konfigurasi ber-$I = +1$, dan
\[
\abs{\{I = +1\}} = 16 \cdot \frac{15!}{2} = \frac{16!}{2}.
\]

\textbf{7.} Langkah yang menggeser ubin di $c$ ke $p$ dibatalkan
dengan menggeser ubin yang sama (kini di $p$) kembali ke $c$: menyusun
$(p\ c)$ dua kali menghasilkan identitas. Karena itu: refleksif (dengan
rangkaian kosong), simetris (balik rangkaiannya, batalkan tiap langkah),
transitif (sambung rangkaiannya) --- sebuah relasi ekuivalensi. Setiap
$\sigma \in R$ memenuhi $I(\sigma) = I(\mathrm{id}) = +1$ menurut
pertanyaan 4, jadi $R \subseteq \{I = +1\}$; dan $\sigma_L \notin R$
memberi kelas kedua.

\textbf{8.} Jika $\sigma(16) = 16$, maka $\sigma$ mempermutasikan
sel $1, \dots, 15$; sebut $\rho$ pembatasan itu. Menambahkan satu
titik tetap tidak mengubah tipe \omterm{def:b2:structures:sn}{siklus} maupun tandanya
(uraikan $\rho$ atas \omterm{def:b2:structures:sn}{transposisi}; hasil kali yang sama berlaku di
$\mathfrak{S}_{16}$), jadi $\varepsilon(\sigma) =
\varepsilon(\rho)$, dan $\chi(16) = +1$ memberi $I(\sigma) =
\varepsilon(\rho)$. Sembarang konfigurasi dapat dibawa ke konfigurasi
yang kekosongannya di rumah: papannya terhubung, jadi jalankan
kekosongan menyusuri lintasan sel bertetangga sampai sel $16$ (tiap
langkahnya sah). Kini andaikan setiap $\rho \in \mathfrak{S}_{15}$ yang
genap terwujud oleh sebuah program. Diberikan $\sigma$ dengan $I(\sigma)
= +1$: jalankan kekosongan pulang ke rumah sehingga tercapai
$\widetilde\sigma$ (yang setara dengan $\sigma$), dengan
$I(\widetilde\sigma) = +1$, yakni pembatasannya $\rho$ genap; program
yang mewujudkan $\rho$ membawa $\widetilde\sigma$ ke
$\widetilde\sigma \circ \rho^{-1} = \mathrm{id}$ (lihat pertanyaan
9). Menurut ketransitifan $\sigma \in R$, sehingga $\{I = +1\} \subseteq
R$ dan keduanya sama.

\textbf{9.} \emph{Satu langkah:} isi $c$ berakhir di $p$
dan kekosongan di $c$, jadi efeknya $\pi = (p\ c)$, dan memang
$\sigma' = \sigma \circ (p\ c) = \sigma \circ \pi^{-1}$.
\emph{Induksi:} jika sebuah rangkaian berefek $\pi_1$ dan membawa
$\sigma$ ke $\sigma \circ \pi_1^{-1}$, maka menyusulinya dengan langkah
berefek $\pi_2 = (p'\ c')$ menghasilkan $(\sigma \circ \pi_1^{-1})
\circ \pi_2^{-1} = \sigma \circ (\pi_2\pi_1)^{-1}$, dan isinya berpindah
menurut $\pi_2 \circ \pi_1$ (mula-mula $\pi_1$, lalu $\pi_2$). Jadi
efeknya tersusun, dan program yang dijalankan dari $\sigma$ berakhir di
$\sigma \circ \pi^{-1}$. \emph{Subgrup:} program kosong berefek
$\mathrm{id}$; penyambungan memberi hasil kali; pembalikan program
(pertanyaan 7) memberi invers. Efek sebuah program membiarkan sel $16$
tetap (kekosongan berawal dan berakhir di rumah), jadi $H \leq
\mathfrak{S}_{15}$. \emph{Kegenapan:} program dengan $k$ langkah punya
$k$ genap (pertanyaan 3), dan $\varepsilon(\sigma \circ \pi^{-1}) =
(-1)^k\varepsilon(\sigma)$ memaksa $\varepsilon(\pi) = +1$: jadi $H
\subseteq A_{15}$.

\textbf{10.} Runut keempat geseran dari kekosongan di $16$: langkah
$16 \to 12$ mengirim isi $12$ ke $16$; langkah $12 \to 11$
mengirim isi $11$ ke $12$; langkah $11 \to 15$ mengirim
isi $15$ ke $11$; langkah $15 \to 16$ mengirim isi
yang terparkir di $16$ (semula di $12$) ke $15$. Hasil bersihnya: $11
\mapsto 12$, $12 \mapsto 15$, $15 \mapsto 11$, kekosongan di rumah, jadi
efeknya $(11\ 12\ 15)$. Perjalanan sebaliknya membatalkannya, dengan
efek $(11\ 12\ 15)^{-1} = (11\ 15\ 12)$. Keduanya efek sebuah program,
jadi keduanya di $H$.

\textbf{11.} Ketetanggaan sel berurutan: di dalam tiap pasangan yang
didaftar, selnya berselisih $1$ pada baris yang sama ($16{-}15$,
$15{-}14$, $14{-}13$; $1{-}2$, $2{-}3$, $3{-}4$; $8{-}7$,
$7{-}6$; $10{-}11$, $11{-}12$) atau berselisih $4$ di dalam satu kolom
($13{-}9$, $9{-}5$, $5{-}1$; $4{-}8$; $6{-}10$; $12{-}16$): jadi ia
jalan tertutup lewat seluruh $16$ sel, sepanjang $16$. Efeknya: seperti
pada pertanyaan 10, dengan menulis sel yang disinggahi $c_0 = 16, c_1 = 15,
\dots, c_{15} = 12$, isi $c_i$ berpindah ke $c_{i-1}$ untuk
$i = 2, \dots, 15$, dan isi $c_1$, yang terparkir di $16$
setelah langkah pertama, terbawa ke $c_{15}$ oleh langkah terakhir.
Jadi efeknya memetakan $15 \mapsto 12$, lalu $14 \mapsto 15$, $13
\mapsto 14$, $9 \mapsto 13$, $5 \mapsto 9$, $1 \mapsto 5$, $2
\mapsto 1$, $3 \mapsto 2$, $4 \mapsto 3$, $8 \mapsto 4$, $7
\mapsto 8$, $6 \mapsto 7$, $10 \mapsto 6$, $11 \mapsto 10$, $12
\mapsto 11$: persis \omterm{def:b2:structures:sn}{siklus} berpanjang $15$ bernama $\zeta$. Urutan siklusnya
berawal $x_0 = 15$, $x_1 = 12$, $x_2 = 11$, dan $(x_0\ x_1\ x_2)
= (15\ 12\ 11)$ memetakan $15 \mapsto 12 \mapsto 11 \mapsto 15$ ---
dan itu persis $(11\ 15\ 12)$, yakni perjalanan dasar yang dibalik.

\textbf{12.} Misalkan $\gamma = (a\ b\ c)$ dan $x \in \intint1n$. Jika
$x = g(a)$: $g\gamma g^{-1}(x) = g(\gamma(a)) = g(b)$; serupa itu
$g(b) \mapsto g(c)$ dan $g(c) \mapsto g(a)$. Jika $x \notin
\{g(a), g(b), g(c)\}$, maka $g^{-1}(x) \notin \{a,b,c\}$ dibiarkan
tetap oleh $\gamma$, sehingga $x$ pun tetap. Jadi $g\gamma g^{-1} =
(g(a)\ g(b)\ g(c))$. Dan untuk $g, h \in H$, berlaku $ghg^{-1} \in H$
menurut aksioma subgrup.

\textbf{13.} Kita punya $\zeta \in H$ (pertanyaan 11) dan $s_0 = (x_0\
x_1\ x_2) \in H$ (pertanyaan 10--11). Karena $\zeta(x_i) = x_{i+1}$
(indeks modulo $15$), pertanyaan 12 memberi
\[
\zeta^{t}\,s_0\,\zeta^{-t}
= \bigl(\zeta^t(x_0)\ \zeta^t(x_1)\ \zeta^t(x_2)\bigr)
= (x_t\ x_{t+1}\ x_{t+2}) = s_t \in H
\qquad (t = 0, 1, \dots, 14).
\]

\textbf{14.} Kecuali pembalikan, andaikan $s = (a\ b\ c)$ dan $t =
(b\ c\ d)$ (\omterm{def:b2:structures:sn}{siklus} berpanjang $3$ pada $\{a,b,c\}$ adalah $(a\ b\ c)$
atau inversnya; demikian pula pada $\{b,c,d\}$; mengganti sebuah
pembangun dengan inversnya tidak mengubah $\langle s, t\rangle$). Maka,
dengan $t$ dikerjakan lebih dulu,
\[
st \colon a \mapsto b,\quad b \mapsto a,\quad c \mapsto d,\quad
d \mapsto c, \qquad\text{yakni}\quad st = (a\ b)(c\ d),
\]
yakni sebuah \omterm{def:b2:structures:sn}{transposisi} ganda. Subgrup $G = \langle s, t\rangle$
terdiri atas permutasi genap atas keempat huruf, jadi $G \leq
A_4$ dan $\abs G \mid 12$; ia memuat unsur berorde $3$
dan unsur berorde $2$, jadi $6 \mid \abs G$ (Lagrange,
\cref{thm:b2:structures:lagrange}, diterapkan pada kedua subgrup
\omterm{def:b2:structures:generated}{siklik} itu). Seandainya $A_4$ punya subgrup $K$ berorde $6$, subgrup itu
berindeks $2$, dan akibatnya $g^2 \in K$ untuk setiap $g \in A_4$: untuk
$g \in K$ hal ini jelas; untuk $g \notin K$ satu-satunya koset adalah
$K$ dan $gK$, jadi koset $g^2K$ sama dengan $K$ atau $gK$, sedangkan
$g^2K = gK$ akan memaksa $g \in K$. Jadi setiap kuadrat terletak di $K$.
Namun setiap \omterm{def:b2:structures:sn}{siklus} berpanjang $3$, sebut saja $\gamma$, adalah kuadrat, sebab
$\gamma = (\gamma^2)^2$, dan $A_4$ memuat delapan \omterm{def:b2:structures:sn}{siklus} berpanjang $3$:
$8 > 6$, kontradiksi. Maka $\abs G = 12$, yakni $G = A_4$.

\textbf{15.} Perpanjang $u \mapsto a$, $v \mapsto b$ menjadi bijeksi
$g_0$ atas $X$ (kirim $k - 2$ huruf sisanya secara bijektif ke mana pun
pada komplemen $\{a, b\}$). Jika $g_0$ ganjil, pilihlah dua huruf yang
berbeda $s_1, t_1 \in X \setminus \{u, v\}$ (mungkin sebab $k \geq 4$)
lalu ganti $g_0$ dengan $g_0 \circ (s_1\ t_1)$, yang genap dan tetap
memetakan $u \mapsto a$, $v \mapsto b$. Perpanjang dengan identitas di
luar $X$: diperoleh permutasi genap $g \in G$ (ia permutasi genap atas
$X$). Lalu menurut pertanyaan 12:
\[
g\,(u\ v\ w)\,g^{-1} = (g(u)\ g(v)\ g(w)) = (a\ b\ w) \in G,
\]
dengan memakai $g(w) = w$.

\textbf{16.} Setiap \omterm{def:b2:structures:sn}{siklus} berpanjang $3$ atas $X \cup \{w\}$ terletak
di $G$: yang tertumpu di $X$ adalah permutasi genap atas $X$; yang
berpenyangga $\{a, b, w\}$ adalah $(a\ b\ w)$ atau $(b\ a\ w)$, dan
keduanya diberikan pertanyaan 15. Menurut \cref{exo:b2:structures:6},
\omterm{def:b2:structures:sn}{siklus} berpanjang $3$ atas himpunan berhuruf $(k+1)$, yaitu $X \cup
\{w\}$, membangun grup alternatingnya, jadi $G$ memuat setiap
permutasi genap atas $X \cup \{w\}$. \emph{Perangkaian:} misalkan $G =
\langle s_0, \dots, s_{12}\rangle$. Lema A yang diterapkan pada $s_0 =
(x_0\ x_1\ x_2)$ dan $s_1 = (x_1\ x_2\ x_3)$ (penyangganya berbagi
$\{x_1, x_2\}$) memberi seluruh permutasi genap atas $X_4 = \{x_0, x_1,
x_2, x_3\}$. Jika $G$ memuat setiap permutasi genap atas $X_m = \{x_0,
\dots, x_{m-1}\}$ ($4 \leq m \leq 14$), maka $s_{m-2} = (x_{m-2}\
x_{m-1}\ x_m)$ punya $u = x_{m-2}, v = x_{m-1} \in X_m$ serta huruf
baru $w = x_m$: Lema B dan bagian pertama tadi memberi seluruh permutasi
genap atas $X_{m+1}$. Induksi sampai $m = 14$ memberi $G \supseteq
A_{15}$ (permutasi genap atas kelima belas sel), sedangkan $G
\subseteq A_{15}$ karena tiap $s_t$ genap: jadi $\langle s_0, \dots,
s_{12}\rangle = A_{15}$.

\textbf{17.} Pertanyaan 13 dan 16 memberi $A_{15} = \langle s_0, \dots,
s_{12}\rangle \subseteq H$; pertanyaan 9 memberi $H \subseteq A_{15}$.
Jadi $H = A_{15}$, berorde $15!/2 = 653\,837\,184\,000$: setiap penataan
ulang yang genap atas kelima belas ubin adalah efek sebuah program.

\textbf{18.} Pertanyaan 8 menyusutkan $R = \{I = +1\}$ menjadi
mewujudkan setiap $\rho \in \mathfrak{S}_{15}$ yang genap oleh sebuah
program, dan itu diselesaikan pertanyaan 17. Bersama pertanyaan 6,
$\abs R = 16!/2 = 10\,461\,394\,944\,000$. \emph{Dua kelas:} biarkan
$t_0 = (14\ 15)$ bekerja pada \emph{isi}: $\varphi(\sigma) = t_0 \circ
\sigma$. Langkah sah dari $\sigma$ juga langkah sah dari
$\varphi(\sigma)$ (sel kekosongannya tak berubah, sebab
$(t_0\sigma)^{-1}(16) = \sigma^{-1}(t_0(16)) = \sigma^{-1}(16)$,
dan sel yang digeser pun sama), lalu $\varphi(\sigma \circ \tau)
= \varphi(\sigma) \circ \tau$: jadi $\varphi$ memetakan rangkaian
langkah ke rangkaian langkah, secara bijektif (ia involusi). Ia
membalik $I$, sebab $\varepsilon(t_0\sigma) = -\varepsilon(\sigma)$
dengan sel kekosongan yang sama. Maka $\varphi$ memetakan kelas $R =
\{I = +1\}$ milik $\mathrm{id}$ secara bijektif pada kelas
$\varphi(\mathrm{id}) = \sigma_L$, yang karenanya seluruhnya $\{I =
-1\}$: jadi tepat ada dua kelas. Inilah teorema Johnson--Story.

\textbf{19.} Indekskan selnya menurut urutan baca dan misalkan $k = 4(i
- 1) + j$ sel kekosongan. Cacah inversi $\sigma$
(pasangan sel $x < y$ dengan $\sigma(x) > \sigma(y)$): pasangan atas
dua sel berubin menyumbang $N$; pasangan yang melibatkan kekosongan:
sel sesudah kekosongan semuanya memuat ubin $< 16$, masing-masing
terbalik (ada $16 - k$ pasangan), sedangkan sel sebelumnya tak pernah
terbalik. Jadi $\varepsilon(\sigma) = (-1)^{N + 16 - k} = (-1)^{N + k}$.
Karena $k = 4(i-1) + j \equiv j \pmod 2$, kita peroleh
\[
I(\sigma) = (-1)^{N + j}\,(-1)^{i + j} = (-1)^{N + i}
= (-1)^{N + r + 1}
\]
dengan memakai $i = 5 - r$. Menurut pertanyaan 18, $\sigma$ terselesaikan
bila dan hanya bila $I(\sigma) = +1$, bila dan hanya bila $N + r$ ganjil.
Periksa: tersusun, $N = 0$, $r = 1$: ganjil, terselesaikan; Loyd, $N =
1$, $r = 1$: genap, tidak terselesaikan.

\textbf{20.} \emph{Aksi:} $e \cdot \sigma = \sigma \circ
\mathrm{id} = \sigma$ dan $g \cdot (h \cdot \sigma) = \sigma
\circ h^{-1} \circ g^{-1} = \sigma \circ (gh)^{-1} = (gh) \cdot
\sigma$; dan $\sigma \circ h^{-1}$ tetap konfigurasi yang
kekosongannya di rumah ($h$ membiarkan sel $16$ tetap). \emph{Bebas:}
dari $\sigma \circ h^{-1} = \sigma$ diperoleh $h^{-1} = \mathrm{id}$
(susun dengan $\sigma^{-1}$). \emph{Orbit = kelas program:} pertanyaan 9
mengatakan konfigurasi yang terjangkau dari $\sigma$ lewat program
persis semua $\sigma \circ \pi^{-1}$, $\pi \in H$, yakni orbit $H
\cdot \sigma$. \emph{Pencacahan:} sifat bebas membuat $h \mapsto h \cdot
\sigma$ injektif, jadi setiap orbit punya $\abs H = 15!/2$ unsur;
karenanya $15!$ konfigurasi berkekosongan di rumah terbelah menjadi
$15!\,/\,(15!/2) = 2$ orbit --- bayangan kedua kelas Johnson--Story pada
konfigurasi berkekosongan di rumah.

\textbf{21.} Papan $3 \times 3$ bersifat bipartit terhadap pewarnaan
papan catur: setiap langkah pada suatu jalan mengubah warnanya, jadi
setiap jalan \emph{tertutup} berpanjang genap. Jalan tertutup yang
menyinggahi masing-masing dari $9$ sel tepat sekali akan berpanjang $9$,
sebuah bilangan ganjil: mustahil. Karena itu konstruksi perjalanan agung
pada Bagian III tidak tersedia untuk teka-teki delapan.

\textbf{22.} \emph{Perjalanan keliling} $9 \to 8 \to 7 \to 4 \to 1
\to 2 \to 3 \to 6 \to 9$ (semua langkahnya bertetangga; berpanjang $8$,
genap): dengan pembukuan pertanyaan 11 memakai $c_1 = 8, c_2 = 7, c_3 =
4, c_4 = 1, c_5 = 2, c_6 = 3, c_7 = 6$, efeknya adalah
\[
\zeta' = (8\ 6\ 3\ 2\ 1\ 4\ 7),
\]
yaitu \omterm{def:b2:structures:sn}{siklus} berpanjang $7$ yang membiarkan pusat $5$ tetap (isi $7$
berpindah ke $8$, isi $4$ ke $7$, isi $1$ ke $4$, isi $2$ ke $1$, isi
$3$ ke $2$, isi $6$ ke $3$, dan isi $8$ ke $6$). \emph{Perjalanan pojok}
$9 \to 6 \to 5 \to 8 \to 9$ berefek $(6\ 8\ 5)$ (isi $5$ berpindah ke
$6$, isi $8$ ke $5$, isi $6$ --- yang terparkir di $9$ --- ke $8$).
Ambil $y_t = \zeta'^{\,t}(8)$: $y_0 = 8, y_1 = 6, y_2 = 3, y_3 = 2,
y_4 = 1, y_5 = 4, y_6 = 7$. Konjugasi (pertanyaan 12) memberi
\[
\zeta'^{\,t}\,(6\ 8\ 5)\,\zeta'^{-t}
= (y_{t+1}\ y_t\ 5) =: T_t \in H_{3\times3},
\]
karena $\zeta'$ membiarkan $5$ tetap. Penyangga $T_0 = (y_1\ y_0\ 5)$
dan $T_1 = (y_2\ y_1\ 5)$ berbagi tepat $\{y_1, 5\}$, jadi Lema A
memberi seluruh permutasi genap atas $\{y_0, y_1, y_2, 5\}$. Lalu
$T_2 = (y_3\ y_2\ 5)$ menambahkan $y_3$ lewat Lema B (kedua hurufnya
$y_2, 5$ berada di himpunan yang sedang dipegang, dengan $k = 4$), dan
$T_3, T_4, T_5$ menambahkan $y_4, y_5, y_6$ berturut-turut: jadi seluruh
permutasi genap atas kedelapan sel bukan-rumah berada di grup program,
yang juga hanya terdiri atas permutasi genap (hujah pertanyaan 9 tidak
bergantung pada bentuk papan). Maka $H_{3\times3} = A_8$, dan penalaran
pertanyaan 6, 8, 18 --- yang juga tak bergantung pada papan ---
menunjukkan konfigurasi yang terjangkau persis yang ber-$I = +1$:
separuh dari $9!$, yakni $181\,440$.

\textbf{23.} Berilah nama sel $0, \dots, n-1$ sepanjang \omterm{def:b2:structures:sn}{siklus} itu.
Satu langkah menukar kekosongan dengan salah satu dari dua tetangganya.
Bacalah ubin itu menurut urutan siklis mulai tepat sesudah kekosongan:
diperoleh kata $w$ yang mendaftar $n - 1$ ubin. Menggeser kekosongan
satu langkah maju mengganti $(p, w)$ dengan $(p + 1, \rho w)$, dengan
$p$ menyatakan sel kekosongan dan $\rho$ memutar kata itu satu posisi;
langkah mundurnya adalah inversnya. Jadi urutan \emph{siklis} ubinnya
(yakni katanya kecuali perputaran) bersifat invarian. Kelas terjangkau
milik $(p, w)$ adalah orbit pemetaan $g \colon (p, w) \mapsto (p+1,
\rho w)$, yaitu unsur berorde $\operatorname{lcm}(n, n-1) =
n(n-1)$ di hasil kali kedua \omterm{def:b2:structures:generated}{grup siklik} itu (translasi
$\Z/n\Z$ dan perputaran atas $n-1$ posisi kata), dengan KPK-nya
sama dengan $n(n-1)$ karena $\gcd(n, n-1) = 1$: jadi tiap kelas punya
tepat $n(n-1)$ konfigurasi, semuanya berkalung sama.
Banyaknya kelas: $n!\,/\,\bigl(n(n-1)\bigr) = (n-2)!$. Untuk $n \geq 5$
berlaku $(n-2)! > 2$: invarian paritas (paling banyak dua kelas) buta
terhadap hampir seluruh halangannya; kekayaan papan $4
\times 4$ --- tempat paritas menjadi \emph{satu-satunya} halangan ---
adalah fakta yang sungguh geometris, bukan fakta formal.

\textbf{24.} Kedua papan memuat ubin dalam urutan yang seluruhnya
terbalik, jadi $N = \binom{15}{2} = 105$ pada kedua kasus (setiap
pasangan ubin terbalik). \emph{Kekosongan di rumah:} $r = 1$, sehingga
$N + r = 106$ genap: tidak terselesaikan. \emph{Kekosongan di sel $1$:}
kekosongan berada di baris teratas, $r = 4$, sehingga $N + r = 109$
ganjil: terselesaikan. Dua papan yang hanya berbeda letak lubangnya
jatuh di sisi tembok yang berlawanan.

\textbf{25.} \emph{Sifat morfisma:} ia mengubah ``satu langkah = satu
\omterm{def:b2:structures:sn}{transposisi}'' menjadi ``satu langkah = satu pembalikan tanda''
(pertanyaan 2 dan 4), sehingga $I$ terhitungkan langkah demi langkah.
\emph{Lagrange:} ia memaksa $6 \mid \abs{\langle s, t\rangle}$ pada Lema
A dan menakar koset pada penyingkiran subgrup berorde $6$ (pertanyaan
14). \emph{Pembangunan oleh \omterm{def:b2:structures:sn}{siklus} berpanjang $3$:} ia mengubah ``$H$
memuat cukup banyak \omterm{def:b2:structures:sn}{siklus} berpanjang $3$'' menjadi ``$H$ memuat seluruh
$A_{15}$'' (pertanyaan 16). \emph{Konjugasi:} ia memproduksi kelima
belas \omterm{def:b2:structures:sn}{siklus} berurutan berpanjang $3$ dari satu perjalanan $2 \times 2$
yang diangkut perjalanan agung (pertanyaan 12--13), dan juga \omterm{def:b2:structures:sn}{siklus}
berpanjang $3$ bertulis $(a\ b\ w)$ pada Lema B. \emph{Asas umumnya:} sebuah invarian
membuktikan kemustahilan, sebuah konstruksi gamblang membuktikan
kemungkinan, dan sebuah masalah tuntas terpecahkan tepat ketika kedua
batas itu bertemu --- di sini, pada seperdua.

\end{solution}
